\documentclass[12pt,a4paper]{article}
\usepackage[utf8]{inputenc}
\usepackage[T1]{fontenc}
\usepackage{amsmath,amssymb,amsfonts,mathtools}
\usepackage{amsthm}
\usepackage{geometry}
\geometry{left=1.0cm,right=1.0cm,top=1.25cm,bottom=3.1cm}
\usepackage{booktabs}
\usepackage{tabularx}
\usepackage{array}
\usepackage{hyperref}
\usepackage{url}
\usepackage{graphicx}
\usepackage{xcolor}
\usepackage{titlesec}
\usepackage{fancyhdr}
\usepackage{microtype}
\usepackage{multicol}
\usepackage{fontspec}
\usepackage{luatexja}          % 日本語組版処理
\usepackage{luatexja-fontspec} % 日本語フォント設定
\usepackage{tikz}
\usepackage{pgfplots}
\pgfplotsset{compat=1.18}
\usetikzlibrary{positioning, arrows.meta, shapes.geometric, calc, decorations.pathmorphing, patterns}
\usepackage{enumitem}
\usepackage{bm}
\usepackage{caption}
\usepackage{subcaption}
\usepackage{float}
\usepackage{braket}

% ========= 日本語フォント設定 (Windows) =========
\defaultfontfeatures{Ligatures=TeX}
\setmainfont{Times New Roman}[
Script=Latin,
Language=English
]
\setsansfont{Meiryo}[
Script=CJK,
Language=Japanese
]
\setmonofont{Courier New}[
Script=Latin,
Language=English
]
% 日本語専用フォント (luatexja-fontspec)
\setmainjfont{Meiryo}[
Script=CJK,
Language=Japanese
]
\setsansjfont{Meiryo}[
Script=CJK,
Language=Japanese
]
\setmonojfont{MS Gothic}[
Script=CJK,
Language=Japanese
]

% ========= YUCT 色 =========
\definecolor{skyblue}{RGB}{0,102,204}
\definecolor{darkblue}{RGB}{0,0,139}
\definecolor{gold}{RGB}{255,215,0}

% ========= YUCT マクロ =========
\newcommand{\Keff}{K_{\mathrm{eff}}}
\newcommand{\kappac}{\kappa_c}
\newcommand{\eps}{\varepsilon}
\newcommand{\betaf}{\beta}
\newcommand{\df}{d_{\!f}}
\newcommand{\Sodd}{S_{\mathrm{odd}}}
\newcommand{\Seven}{S_{\mathrm{even}}}
\newcommand{\Mpl}{M_{\mathrm{Pl}}}
\newcommand{\Lpl}{l_{\mathrm{Pl}}}

% ========= 定理環境 (日本語) =========
\newtheorem{theorem}{定理}[section]
\newtheorem{lemma}[theorem]{補題}
\newtheorem{proposition}[theorem]{命題}
\newtheorem{definition}[theorem]{定義}
\newtheorem{remark}[theorem]{注記}
\renewcommand{\proofname}{証明}

\DeclareMathOperator{\Ent}{Ent}
\DeclareMathOperator{\Tr}{Tr}

% YUCT コマンド
\newcommand{\Dir}{\mathcal{D}}
\newcommand{\cL}{\mathcal{L}}
\newcommand{\Planck}{\mathrm{Pl}}

% ========= セクション書式 =========
\titleformat{\section}{\LARGE\bfseries\color{darkblue}}{\thesection}{1em}{}
\titleformat{\subsection}{\Large\bfseries\color{darkblue}}{\thesubsection}{1em}{}

% ========= ヘッダ・フッタ =========
\fancypagestyle{firstpage}{%
	\fancyhf{}%
	\renewcommand{\headrulewidth}{0pt}%
	\renewcommand{\footrulewidth}{0pt}%
	\fancyfoot[C]{%
		\begin{minipage}[t]{\textwidth}
			\centering
			\textcolor{gold}{\rule{\textwidth}{1.6pt}}\\[5pt]
			{\small \textsuperscript{1}Yakushev Research, \url{https://yuct.org/}} \hfill
			{\small YPSDCプロトコル: \url{https://ypsdc.com/}}\\[-2pt]
			\textcolor{gold}{\rule{\textwidth}{1.6pt}}\\[5pt]
			\textbf{ヤクシェフ統一調整理論 2026} \hfill \thepage\\[10pt]
		\end{minipage}%
	}%
}

\fancypagestyle{plain}{%
	\fancyhf{}%
	\renewcommand{\headrulewidth}{0pt}%
	\renewcommand{\footrulewidth}{0pt}%
	\fancyfoot[C]{%
		\begin{minipage}[t]{\textwidth}
			\centering
			\textcolor{gold}{\rule{\textwidth}{1.6pt}}\\[4pt]
			\textbf{ヤクシェフ統一調整理論 2026} \hfill \thepage\\[4pt]
		\end{minipage}%
	}%
}

\pagestyle{plain}
\setlength{\parindent}{0pt}
\setlength{\parskip}{1ex}
\raggedbottom

\hypersetup{
	colorlinks=true,
	linkcolor=blue,
	citecolor=blue,
	urlcolor=blue,
}

% ========= ヘッダマクロ =========
\newcommand{\makeheader}{%
	\noindent
	\begin{minipage}{0.17\textwidth}
		\raggedright
		{\sffamily\bfseries\color{skyblue}\fontsize{46}{46}\selectfont YUCT}%
	\end{minipage}%
	\hspace{30pt}
	\begin{minipage}{0.32\textwidth}
		\raggedright
		{\sffamily\fontsize{11}{13}\selectfont ヤクシェフ\\ 統一\\ 調整理論}%
	\end{minipage}%
	\hfill
	\begin{minipage}{0.38\textwidth}
		\raggedleft
		{\sffamily\bfseries 技術ノート}\\ 
		{\sffamily 2026年5月発行}\\ 
		{\sffamily\footnotesize \scriptsize\url{https://doi.org/10.5281/zenodo.18444598}}%
	\end{minipage}
	\par\vspace{5pt}
	\noindent\textcolor{gold}{\rule{\textwidth}{1.6pt}}
	\par\vspace{0pt}
}

\begin{document}
	\thispagestyle{firstpage}
	\makeheader
	
	\vspace{0cm}
	{\LARGE\bfseries 調整ネットワークの幾何学的張力としての重力と\\
		YUCT における循環宇宙論 \\
		（厳密な導出を含む最終版）\par}
	\vspace{0.2cm}
	
	{\large Alexey V. Yakushev\textsuperscript{1}} \\
	\vspace{0.0cm}
	
	{\color{darkblue}\textbf{\LARGE 要旨}} \par
	{\large
		\noindent
		我々は、ヤクシェフ統一調整理論 (YUCT) の枠組みにおいて、重力を創発現象として数学的に厳密に導出する。調整効率 \(K_{\mathrm{eff}}\) の微視的定義とスカラー場 \(\phi = \ln (K_{\mathrm{eff}}/K_0)\) から出発し、調整ネットワークの自由エネルギー汎関数を構成する。すべての定数は先験的辞書のパラメータを通じて表現される。作用の変分は、暗黒物質と暗黒エネルギーの両方を置き換える幾何学的張力テンソルを持つ修正アインシュタイン方程式を与える。ニュートン極限は正しく導出され、追加項が平坦な回転曲線を生成する有効暗黒物質密度として機能することを示す。宇宙定数 \(\Omega_\Lambda = 2/3\) は 19 次元予調整空間のトポロジーから導かれ、\(K_{\mathrm{eff}} = \Phi\)（黄金比）での相転移から循環宇宙論が現れる。予測は銀河および宇宙論的スケールで定量化される。}
	
	\vspace{0.1cm}
	\noindent
	{\color{darkblue}\textbf{キーワード:}} YUCT，調整効率 \(K_{\mathrm{eff}}\)，幾何学的張力，修正重力，暗黒物質，暗黒エネルギー，宇宙定数 \(\Omega_\Lambda=2/3\)，循環宇宙論，黄金比，トポロジカル不変量，19 次元予調整空間。
	
	\vspace{0.4cm}
	
	\tableofcontents
	
	\section{調整場の微視的基礎}
	
	\subsection{調整クラスタと効率}
	\begin{definition}
		レベル \(n\) の調整クラスタは \(\mathcal{C}_n = (\mathcal{O}_n, \Dir_n, L_n, \tau_n, K_{\mathrm{eff},n})\) で定義される。ここで \(\mathcal{O}_n\) はエージェント集合，\(\Dir_n\) は先験的辞書，\(L_n\) はサイズ，\(\tau_n\) は調整時間であり，
		\[
		K_{\mathrm{eff},n} = \frac{H(\mathcal{A})}{H(\mathcal{K})}
		\]
		は作用エントロピーと索引エントロピーの比である。
	\end{definition}
	
	\(N\) 個のエージェントからなるクラスタに対して、有効辞書サイズは \(|\Dir_{\mathrm{eff}}| \propto \mathcal{N}^{\alpha N}\) と成長する。ここで \(\mathcal{N}\) は基本アルファベットサイズ，\(\alpha<1\) は相関指数である。辞書エントロピーは
	\[
	H(\Dir) = \log_2 |\Dir_{\mathrm{eff}}| = \alpha N \log_2 \mathcal{N}.
	\]
	各索引が 1 ビットを運ぶと仮定すると \(H(\mathcal{K}) = N\) であり、したがって
	\begin{equation}
		K_{\mathrm{eff}} = \alpha \log_2 \mathcal{N}. \label{eq:Keff-micro}
	\end{equation}
	単一クラスタでは \(\alpha\) と \(\mathcal{N}\) は固定され、\(K_{\mathrm{eff}}\) は定数である。連続極限では局所的な深さ（したがって \(\alpha\)）と局所的な豊かさ（\(\mathcal{N}\)）は空間的にゆっくり変化しうる。支配的な空間依存性は \(\mathcal{N}\) を通じて現れ、\(\alpha\) の変動は副次的であり参照スケールの再正規化に吸収される。そこでスカラー場を
	\begin{equation}
		\phi(x) = \ln \frac{K_{\mathrm{eff}}(x)}{K_0}, \qquad K_0 \equiv \alpha_0 \log_2 \mathcal{N}_0,
		\label{eq:phi-def}
	\end{equation}
	と定義する。ここで \(\alpha_0,\mathcal{N}_0\) は真空（プランクスケール）配置を指す。プランクスケールではチャネル容量は \(1/t_P\) であり辞書は最小：\(\mathcal{N}_0 \approx 2\)，\(\alpha_0 \rightarrow 1\)，したがって \(K_0 \approx 1\) かつ真空中で \(\phi=0\) である。
	
	\subsection{辞書のエントロピー密度}
	辞書エントロピーの空間密度は \(s = H(\Dir)/V_c\) である。クラスタの相関体積は \(V_c \propto N^{1/d_f}\) でスケールする。ここで \(d_f\) はフラクタル次元である。\(N \propto K_{\mathrm{eff}}^{1/\alpha}\) と (\ref{eq:Keff-micro}) を用いると、大きなクラスタに対して
	\begin{equation}
		s(\phi) = s_0\, e^{\lambda\phi}, \qquad \lambda = \frac{1}{\beta} = \frac{3}{2},
		\label{eq:s-phi}
	\end{equation}
	が得られる。ここで \(\beta = 2/3\) は普遍誤差指数である（付録 L）。定数 \(s_0\) は真空エントロピー密度である。
	
	\section{調整ネットワークの自由エネルギー}
	
	\subsection{ヘルムホルツ汎関数}
	調整ネットワークはエントロピーを運ぶ弾性媒体である。その平衡はヘルムホルツ自由エネルギーの最小化から得られる：
	\begin{equation}
		F = \int d^4x \sqrt{-g} \left[ \frac{1}{2}\kappa (\nabla\phi)^2 + V(\phi) - T s(\phi) \right],
		\label{eq:F}
	\end{equation}
	ここで \(T\) はネットワークの有効温度である。項 \(-Ts\) は標準的な \(F = U - TS\) に対応する。ここで \(U = \int [\frac12\kappa(\nabla\phi)^2 + V(\phi)]\) である。
	
	\subsection{ポテンシャルの決定}
	真空 \(\phi=0\) は極値でなければならない：\(\delta F/\delta\phi|_0 = 0\)。\(s'(\phi) = \lambda s_0 e^{\lambda\phi}\) を用いると
	\[
	V'(0) = T s'(0) = \lambda T s_0.
	\]
	この条件を満たし、大きな \(\phi\) でのエントロピー成長と整合する最も単純なポテンシャルは
	\begin{equation}
		\boxed{V(\phi) = V_0 \left( e^{\lambda\phi} - \lambda\phi - 1 \right), \qquad V_0 \equiv T s_0.}
		\label{eq:V}
	\end{equation}
	\(\phi \ll 1\) では \(V(\phi) \approx \frac12 V_0 \lambda^2 \phi^2\)；\(\phi \gg 1\) では \(V(\phi) \sim V_0 e^{\lambda\phi}\) となり自由エネルギー中のエントロピー項を打ち消し、\(\phi=0\) が唯一の大域的極小であることを保証する。
	
	\subsection{\(\kappa\) と \(V_0\) の微視的推定}
	剛性 \(\kappa\) は辞書効率の単位勾配あたりのエネルギーコストである。次のように推定できる：
	\[
	\kappa \approx \frac{E_P}{\ell_P} \cdot \frac{1}{\alpha_0 \log_2 \mathcal{N}_0} \sim \frac{\hbar c}{\ell_P^2},
	\]
	ここで \(E_P, \ell_P\) はプランクエネルギーとプランク長である。真空エントロピー密度はプランクセルあたり約 1 ビット：\(s_0 \approx k_B / \ell_P^3\)。\(T \sim T_P = \sqrt{\hbar c^5 / G k_B^2}\) を取ると
	\[
	V_0 = T s_0 \sim \frac{\hbar c}{\ell_P^4} \sim M_{\mathrm{Pl}}^4.
	\]
	これらの粗い推定はパラメータをプランクスケールに固定するが、その精密な値は低エネルギー現象（銀河回転と宇宙膨張）への適合によって決定される。これは以下で行う。
	
	\section{修正アインシュタイン方程式}
	
	調整場を重力に結合させると全作用は
	\begin{equation}
		S = \int d^4x \sqrt{-g} \left[ \frac{1}{16\pi G_0} R + \frac{1}{2}\kappa (\nabla\phi)^2 + V(\phi) \right].
		\label{eq:action}
	\end{equation}
	\(g^{\mu\nu}\) で変分すると
	\begin{equation}
		G_{\mu\nu} = 8\pi G_0 \left( \Theta_{\mu\nu} + T_{\mu\nu}^{(\mathrm{matter})} \right),
		\label{eq:einstein}
	\end{equation}
	ここで \textbf{幾何学的張力テンソル} は
	\begin{equation}
		\Theta_{\mu\nu} = \kappa \nabla_\mu \phi \nabla_\nu \phi - g_{\mu\nu} \left[ \frac{\kappa}{2} (\nabla\phi)^2 + V(\phi) \right].
		\label{eq:Theta}
	\end{equation}
	\(\phi\) による変分は運動方程式を与える：
	\begin{equation}
		\kappa \Box \phi - V'(\phi) = 0, \qquad V'(\phi) = \lambda V_0 \left( e^{\lambda\phi} - 1 \right).
		\label{eq:phieom}
	\end{equation}
	
	\section{ニュートン極限と有効暗黒物質}
	
	\subsection{\(\phi\) の線形化場方程式}
	静的球対称背景計量
	\[
	ds^2 = -(1+2\Phi)dt^2 + (1-2\Phi)\delta_{ij}dx^i dx^j,
	\]
	を考え、\(\phi = \phi_{\mathrm{vac}} + \delta\phi(r)\) と書く。ここで \(\phi_{\mathrm{vac}} = 0\)。\(|\delta\phi| \ll 1\) に対して \(V'(\delta\phi) \approx \lambda^2 V_0 \delta\phi\) を展開し、(\ref{eq:phieom}) から
	\begin{equation}
		\kappa \nabla^2 \delta\phi - \lambda^2 V_0 \delta\phi = 0.
		\label{eq:lin-phi}
	\end{equation}
	無限遠でゼロとなる解は
	\begin{equation}
		\delta\phi(r) = \frac{A}{r} e^{-r/\ell}, \qquad \ell \equiv \sqrt{\frac{\kappa}{\lambda^2 V_0}}.
		\label{eq:deltaphi-sol}
	\end{equation}
	
	\subsection{ポアソン方程式}
	弱場極限において、エネルギー-運動量テンソル (\ref{eq:Theta}) を組み合わせたアインシュタイン方程式は次の修正ポアソン方程式を導く（詳細な導出は付録 G）：
	\begin{equation}
		\nabla^2 \Phi = 4\pi G_0 \rho_b + \frac{\kappa}{2} \nabla^2 \delta\phi.
		\label{eq:poisson}
	\end{equation}
	第２項は有効的な追加密度として働く：
	\begin{equation}
		\rho_{\mathrm{DM}}^{\mathrm{eff}} \equiv \frac{\kappa}{8\pi G_0} \nabla^2 \delta\phi.
		\label{eq:rhoDM}
	\end{equation}
	この結果は \(\delta\phi\) の線形オーダーで正確であり、勾配エネルギー \(\frac{\kappa}{2}(\nabla\delta\phi)^2\) は２次オーダーに寄与し、線形化ポアソン方程式には影響しない。
	
	\subsection{銀河回転曲線}
	銀河スケール (\(r \ll \ell\)) に対して、スカラー場の輪郭は \(\delta\phi(r) = A/r\) であり、ここで \(A\) は長さの次元を持つ定数である。これを修正ポアソン方程式 (\ref{eq:poisson}) に代入し、一度積分すると動径加速度が得られる：
	\[
	\frac{d\Phi}{dr} = \frac{G_0 M_b(r)}{r^2} + \frac{\kappa}{2r^2} \int_0^r \nabla^2\delta\phi \, r'^2 dr'.
	\]
	\(\nabla^2(1/r) = -4\pi\delta^{(3)}(\mathbf{r})\) を用いると、任意の \(r>0\) に対して積分は \(-A\) に等しい。したがって
	\[
	\frac{d\Phi}{dr} = \frac{G_0 M_b(r)}{r^2} - \frac{\kappa A}{2r^2}.
	\]
	拡張された質量分布に対して、重子項 \(G_0 M_b(r)/r^2\) は \(1/r^2\) よりも遅く減少する；組み合わせ \(M_b(r) - \frac{\kappa A}{2G_0}\) は大 \(r\) で定数に近づく（重子質量が飽和するため）。その結果、円周速度
	\[
	v_c^2 = r\frac{d\Phi}{dr} = \frac{G_0 M_b(r)}{r} - \frac{\kappa A}{2r}
	\]
	は定数に近づく。物理的理由は、有効暗黒物質密度 \(\rho_{\mathrm{DM}}^{\mathrm{eff}} = \frac{\kappa}{8\pi G_0} \nabla^2\delta\phi\) が重子コアの外側で \(r^{-2}\) として振る舞い、累積質量 \(M_{\mathrm{DM}}(r) \propto r\) と \(r\) に依存しない追加の加速度を生じるためである。
	
	漸近速度を基本パラメータで表すために、無次元パラメータ
	\begin{equation}
		\alpha \equiv \frac{\kappa A}{G_0 M_b^2},
		\label{eq:alpha}
	\end{equation}
	を定義する。これは重力に対する調整張力の強さを特徴づける。漸近円周速度の二乗は
	\begin{equation}
		v_c^2(\infty) = \alpha\, G_0 M_b.
		\label{eq:vc-squared}
	\end{equation}
	観測は \(M_b \sim 10^{11}M_\odot\) の銀河に対して \(v_c \approx 200\ \mathrm{km/s}\) を与える；\(G_0 = 6.67\times 10^{-11}\ \mathrm{m^3\,kg^{-1}\,s^{-2}}\) と \(M_b \sim 2\times 10^{41}\ \mathrm{kg}\) を用いると、\(\alpha \approx 2 \times 10^{-7}\)（SI単位）を得る。自然単位系 \(G_0 = 1/M_{\mathrm{Pl}}^2\) では、この \(\alpha\) はオーダー１であり、調整張力が銀河スケールで重力と同程度の強さであることを反映する。積 \(\kappa A\) は \(\alpha\) と銀河の重子質量によって完全に決定される。
	
	コンプトン波長 \(\ell = \sqrt{\kappa/(\lambda^2 V_0)}\) は条件 \(r_{\mathrm{gal}} \ll \ell\)（\(1/r\) 近似を使うため）によって制約され、\(\ell \sim \text{数}\times 10\ \mathrm{kpc}\) を与える。これは典型的な暗黒物質ハローのサイズと一致する。
	
	\section{宇宙定数としてのトポロジカル不変量}
	
	完全な 19 次元 YUCT（付録 A）では、予調整場は 18 次元のファイバー \(F\) を持つ多様体上に生きる。\(F\) 上の予調整演算子のトポロジカル指数は、真空エネルギーに Casimir 効果を通じて寄与する独立な調和形式を正確に 12 個与える。結果として、
	\begin{equation}
		\Omega_\Lambda = \frac{12}{18} = \frac{2}{3}.
		\label{eq:OmegaLambda}
	\end{equation}
	この比は \(V(\phi)\) の詳細に依存しない。完全な導出は付録 A に示す。
	
	一様等方性背景において、場 \(\phi\) の時間発展は
	\[
	\ddot\phi + 3H \dot\phi + \frac{\lambda V_0}{\kappa}(e^{\lambda\phi}-1) = 0.
	\]
	数値解は、\(\phi\) が状態方程式 \(w_\phi\) が \(-1\) に近づくトラッキング領域に引き寄せられ、現在のエネルギー密度分数が値 \(2/3\) に達することを示している。したがって、ポテンシャルパラメータは現在の膨張率と、そのアトラクターが今日に達するという条件によって制約される。
	
	\subsection{トポロジカル不変量からの微細構造定数}
	
	宇宙定数を固定する同じ整数 \(N_{\mathrm{gauge}} = 12\) が、電磁相互作用の強度も支配する。  
	YUCT では、ゲージ場は調整ネットワークの無質量励起として現れ、その低エネルギー有効結合は緊密ファイバー \(F_{18}\) 上の独立な調和モードの数によって決定される。  
	プランクスケールでは、12 個のモードすべてが活性であり、逆結合定数は
	\begin{equation}
		\alpha_0^{-1} = \left( \frac{S_{\mathrm{odd}}}{S_{\mathrm{even}}} \right)^{12}
		= \left( \frac{3}{2} \right)^{12} \approx 129.746 .
	\end{equation}
	
	宇宙が冷却して電弱対称性が破れると、大質量の \(W\) および \(Z\) ボソンが脱結合し、有効な寄与モード数は小さな普遍的修正を受ける。  
	この修正は、重力の温度依存性に現れる同じパラメータ \(\beta\gamma\)（付録 P）と基本比 \(S_{\mathrm{even}}/S_{\mathrm{odd}}\) の積に比例する：
	\begin{equation}
		\delta N = \beta\gamma \, \frac{S_{\mathrm{even}}}{S_{\mathrm{odd}}} .
		\label{eq:deltaN}
	\end{equation}
	経験値 \(\beta\gamma = 0.20\)（付録 P）と \(S_{\mathrm{even}}/S_{\mathrm{odd}} = 2/3\) を用いると、\(\delta N \approx 0.1333\) を得る。
	
	したがって、実験室スケールでの電磁結合を決定する有効なゲージ自由度は
	\begin{equation}
		N_{\mathrm{eff}} = 12 + \delta N = 12 + \beta\gamma\,\frac{S_{\mathrm{even}}}{S_{\mathrm{odd}}} \approx 12.1333 .
	\end{equation}
	予測される逆微細構造定数は
	\begin{equation}
		\boxed{\alpha^{-1}_{\mathrm{YUCT}} = \left( \frac{S_{\mathrm{odd}}}{S_{\mathrm{even}}} \right)^{N_{\mathrm{eff}}}
			= \left( \frac{3}{2} \right)^{12 + \beta\gamma\,S_{\mathrm{even}}/S_{\mathrm{odd}}} } .
		\label{eq:alphaYUCT}
	\end{equation}
	数値的に、
	\[
	\alpha^{-1}_{\mathrm{YUCT}} = \left( \frac{3}{2} \right)^{12.1333} \approx 137.0 ,
	\]
	これは CODATA 2022 値 \(\alpha^{-1}_{\mathrm{exp}} = 137.035999084\) と \(0.03\%\) 未満の差しかない。
	
	式 (\ref{eq:alphaYUCT}) には自由パラメータは含まれない：\(12\) は 19 次元調整空間のトポロジカル不変量、\(S_{\mathrm{odd}}/S_{\mathrm{even}}\) と \(\beta\gamma\) は独立な観測によってすでに固定された YUCT の基本定数である（付録 Ferra1.2 および P）。  
	この結果は、伝統的に標準模型の自由入力パラメータであった微細構造定数が、実際には調整フレームワーク内で導出される量であることを示している。
	
	\subsection{フラクタル梯子：結合定数と質量階層性のトポロジカル起源}
	\label{sec:fractal_ladder}
	
	整数 \(N_{\mathrm{gauge}}=12\) と基本比 \(S_{\mathrm{odd}}/S_{\mathrm{even}}=3/2\) から微細構造定数を導出することは孤立した結果ではない。それは、素粒子から生きた細胞に至るすべての安定な物体の質量を生成する普遍的なフラクタル梯子の低エネルギーアンカーである。
	
	図~\ref{fig:fractal_ladder} はこの梯子を視覚化したものである。左軸は半整数指数 \(N_f\) で、これは質量スペクトル \(m = m_e q^{N_f}\) を量子化し、\(q = (3/2)^{1/3}\) である。右軸は有効調整深度 \(L = N_f/3\) である。梯子の各段は \(N_f\) の整数または半整数値に対応する。
	
	\begin{figure}[htbp]
		\centering
		\begin{tikzpicture}[scale=0.95]
			% ===== 左軸：質量梯子（対数スケール）=====
			\draw[->] (0,0) -- (0,14) node[above,align=center] {\(N_f\) \\ （半整数指数）};
			\draw[->] (0,0) -- (15,0) node[right] {物体};
			% 目盛りとラベル
			\foreach \y/\lab in {0/{$0$ (e)}, 10/{$10$ ($u$)}, 20/{}, 30/{}, 40/{$\mu$}, 50/{}, 60/{$\tau$}, 70/{}, 80/{}, 90/{$t$}, 100/{}, 110/{}, 120/{ユビキチン}, 130/{}, 140/{ヌクレオソーム}, 150/{}, 160/{リボソーム}, 170/{}, 180/{}, 190/{NPC}, 200/{}, 210/{}, 220/{}, 230/{}, 240/{}, 250/{}, 260/{大腸菌}, 270/{}, 280/{}, 290/{}, 300/{HeLa細胞}, 310/{}} {
				\draw[gray, thin] (0,\y/24) -- (15,\y/24);
				\node[left,font=\tiny] at (0,\y/24) {\lab};
			}
			% ===== 右軸：深度 L = N_f/3 =====
			\draw[->] (15.5,0) -- (15.5,14) node[above,align=center] {\(L\) \\ （調整深度）};
			\foreach \y/\lab in {0/0, 3/1, 6/2, 9/3, 12/4, 15/5, 18/6, 21/7, 24/8, 27/9, 30/10, 33/11, 36/12, 39/13, 42/14, 45/15, 48/16, 51/17, 54/18, 57/19, 60/20} {
				\node[right,font=\tiny] at (15.5,\y/4) {\lab};
			}
			% ===== トップセクション：トポロジーボックス =====
			\draw[fill=blue!10, rounded corners] (0.5,12.5) rectangle (15,14);
			\node[font=\bf,align=center] at (7.75,13.25) {19 次元幾何学：\( \mathcal{M}_{19} = M_4 \times F_{18} \)};
			\node[font=\small,align=center] at (7.75,12.75) {\(F_{18}\) 上の \(N_{\mathrm{gauge}} = 12\) 個の調和モード};
			% トポロジーからゲージ結合への矢印
			\draw[->, thick, blue] (7.75,12.5) -- (7.75,11.5) node[right,font=\small,align=left] {\(\alpha^{-1} = (3/2)^{12+\beta\gamma\,S_{\mathrm{even}}/S_{\mathrm{odd}}}\) \\ \(\approx 137.036\)（微細構造定数）};
			\draw[->, thick, blue] (7.75,12.5) -- (2,11.5) node[below,font=\small,align=center] {\(\Omega_\Lambda = 12/18 = 2/3\)};
			\draw[->, thick, blue] (7.75,12.5) -- (13,11.5) node[below,font=\small,align=center] {\(\alpha_s^{-1} = (3/2)^{8+\beta\gamma\,S_{\mathrm{even}}/S_{\mathrm{odd}}}\)};
			% ===== 質量梯子：素粒子 =====
			\draw[fill=red!20, rounded corners] (0.5,9.5) rectangle (15,11);
			\node[font=\bf] at (7.75,10.5) {基本フェルミオン：\(m_f = m_e q^{N_f}\)};
			\node[font=\small] at (7.75,10.0) {\(N_f \in \frac12\mathbb{Z}\)，偏差 \(< 1\%\)};
			% ===== 質量梯子：原子核 =====
			\draw[fill=blue!20, rounded corners] (0.5,8.0) rectangle (15,9.2);
			\node[font=\bf] at (7.75,8.85) {原子核};
			\node[font=\small] at (7.75,8.35) {陽子，重陽子，\(^4\)He，\(^{12}\)C，\(^{16}\)O……};
			% ===== 質量梯子：タンパク質複合体 =====
			\draw[fill=green!20, rounded corners] (0.5,6.5) rectangle (15,7.7);
			\node[font=\bf] at (7.75,7.35) {タンパク質複合体};
			\node[font=\small] at (7.75,6.85) {ユビキチン，ヘモグロビン，リボソーム，プロテアソーム，NPC};
			% ===== 質量梯子：生きた細胞 =====
			\draw[fill=orange!20, rounded corners] (0.5,5.0) rectangle (15,6.2);
			\node[font=\bf] at (7.75,5.85) {生きた細胞};
			\node[font=\small] at (7.75,5.35) {細菌（大腸菌），HeLa細胞，線維芽細胞};
			% ===== 宇宙論スケール =====
			\draw[fill=purple!20, rounded corners] (0.5,3.5) rectangle (15,4.7);
			\node[font=\bf] at (7.75,4.35) {宇宙論スケール};
			\node[font=\small] at (7.75,3.85) {\(M_{\mathrm{bar}}/m_p = (M_{\mathrm{Pl}}/m_e)^{3.5}\)};
			% ===== ボトムセクション：普遍定数 =====
			\draw[fill=yellow!20, rounded corners] (0.5,1.5) rectangle (15,3.0);
			\node[font=\bf,align=center] at (7.75,2.5) {基本量子：\(q = (3/2)^{1/3} \approx 1.1447\)};
			\node[font=\small,align=center] at (7.75,2.0) {\(S_{\mathrm{odd}} = 6/5 = 1.2\)，\(S_{\mathrm{even}} = 4/5 = 0.8\)，\(\beta = 2/3\)，\(\sigma = 0.20\)};
			% ===== 接続を示す中央の矢印 =====
			\draw[->, thick, black!70, decorate, decoration={snake, amplitude=0.3mm, segment length=2mm}] (14,11) -- (14,1.5) node[midway,right,font=\small,align=left] {普遍的 \\ 質量梯子 \\ \(m = m_e q^{N_f}\)};
		\end{tikzpicture}
		\caption{\textbf{物理的実在のフラクタル梯子。} 緊密ファイバー \(F_{18}\) 上のトポロジカル整数 \(N_{\mathrm{gauge}}=12\) はゲージ結合を固定する（上部ボックス）。同じ基本比 \(3/2 = S_{\mathrm{odd}}/S_{\mathrm{even}}\) が普遍的質量梯子 \(m = m_e q^{N_f}\) を生成する。ここで \(q = (3/2)^{1/3}\) であり、基本フェルミオンから原子核、タンパク質複合体、生きた細胞、そして宇宙論スケールまで広がる。修正 \(\beta\gamma\,S_{\mathrm{even}}/S_{\mathrm{odd}} \approx 0.1333\) は有効モード数を 12 から 12.1333 に移し、\(\alpha^{-1} \approx 137.036\) を与える。}
		\label{fig:fractal_ladder}
	\end{figure}
	
	\subsubsection{梯子の三つの段}
	
	フラクタル梯子は三つの相互に連結した段に基づいており、それぞれが同じ基本比 \(S_{\mathrm{odd}}/S_{\mathrm{even}} = 3/2\) によって支配される：
	
	\begin{enumerate}[leftmargin=*]
		\item \textbf{ゲージ結合（上部）。} 整数 \(N_{\mathrm{gauge}} = 12\) は 19 次元調整空間のトポロジカル不変量である。逆微細構造定数は \(\alpha^{-1} = (3/2)^{12.1333} \approx 137.036\) であり、ここで小さな修正 \(\beta\gamma\,S_{\mathrm{even}}/S_{\mathrm{odd}} \approx 0.1333\) は電弱スケール以下での重ボソンの脱結合を考慮している。これは梯子の \emph{最低段} である：電磁相互作用の強度は緊密ファイバー上の調和モードの数によって固定される。
		\item \textbf{フェルミオン質量（中部）。} スケーリング量子 \(q = (3/2)^{1/3}\) と半整数量子化 \(m_f = m_e q^{N_f}\) は、すべての荷電フェルミオンの質量をサブパーセント精度で生成する。同じ量子はハドロン、原子核、タンパク質複合体の質量も支配する。
		\item \textbf{宇宙定数（上部）。} 有効ゲージモード数とファイバーの全次元の比は \(\Omega_\Lambda = 12/18 = 2/3\) を与え、連続パラメータなしで暗黒エネルギー密度を固定する。
	\end{enumerate}
	
	この梯子は、\textbf{フレームワーク全体に自由連続パラメータが存在しない} ことを示している。整数 12、比 \(3/2\)、および修正 \(\beta\gamma\,S_{\mathrm{even}}/S_{\mathrm{odd}}\) はすべて、19 次元調整空間のトポロジーと YUCT の代数的ループによって固定されている。したがって、電子の質量、電磁相互作用の強さ、そして宇宙の幾何学は、すべて同じフラクタルパターンの異なる投影なのである。
	
	\subsection{調整深度からのゲージ結合と湯川結合}
	
	\(\Omega_\Lambda\) を固定する整数 \(N_{\mathrm{gauge}} = 12\) は電弱相互作用と強い相互作用の強度も支配し、フェルミオン–ヒッグス湯川結合は関与する粒子の調整深度から直接導かれる。
	
	\subsubsection{微細構造定数}
	
	プランクスケールでは、12 個のゲージモードすべてが活性であり、逆電磁結合は
	\begin{equation}
		\alpha_0^{-1} = \left(\frac{S_{\mathrm{odd}}}{S_{\mathrm{even}}}\right)^{\!12}
		= \left(\frac{3}{2}\right)^{\!12} \approx 129.746 .
	\end{equation}
	電弱スケール以下では、大質量の \(W\) および \(Z\) ボソンが脱結合し、有効な寄与モード数は積 \(\beta\gamma\)（付録 P）と基本比 \(S_{\mathrm{even}}/S_{\mathrm{odd}}\) に比例した普遍的な修正だけ減少する：
	\begin{equation}
		\delta N = \beta\gamma\,\frac{S_{\mathrm{even}}}{S_{\mathrm{odd}}}\, .
		\label{eq:deltaN2}
	\end{equation}
	\(\beta\gamma = 0.20\) と \(S_{\mathrm{even}}/S_{\mathrm{odd}} = 2/3\) を用いると \(\delta N \approx 0.1333\) を得る。したがって実験室スケールでの有効ゲージ自由度は
	\begin{equation}
		N_{\mathrm{eff}} = 12 + \delta N = 12 + \beta\gamma\,\frac{S_{\mathrm{even}}}{S_{\mathrm{odd}}}
		\approx 12.1333 .
	\end{equation}
	予測される逆微細構造定数は
	\begin{equation}
		\boxed{\alpha^{-1}_{\mathrm{YUCT}} =
			\left(\frac{S_{\mathrm{odd}}}{S_{\mathrm{even}}}\right)^{\!N_{\mathrm{eff}}}
			= \left(\frac{3}{2}\right)^{\!12 + \beta\gamma\,S_{\mathrm{even}}/S_{\mathrm{odd}}} } .
		\label{eq:alphaYUCT2}
	\end{equation}
	数値的に、
	\[
	\alpha^{-1}_{\mathrm{YUCT}} = \left(\frac{3}{2}\right)^{\!12.1333} \approx 137.0 ,
	\]
	これは CODATA 2022 値 \(\alpha^{-1}_{\mathrm{exp}} = 137.035999084\) と \(0.03\%\) 未満の差しかない。式 (\ref{eq:alphaYUCT2}) には自由パラメータは含まれない：\(12\) は 19 次元調整空間のトポロジカル不変量、\(S_{\mathrm{odd}}/S_{\mathrm{even}}\) と \(\beta\gamma\) は独立な観測ですでに固定された YUCT の基本定数である（付録 Ferra1.2 および P）。
	
	\subsubsection{弱い結合}
	弱い相互作用は別個の有効モード数を必要としない。ワインバーグ角 \(\theta_W\) は質量比 \(m_W/m_Z = \cos\theta_W\) によって決定され、\(m_W\) と \(m_Z\) はどちらも深度 \(L_W = 96\) と適切な共鳴因子から導かれる（付録 \(\Lambda\)）。したがって弱結合 \(\alpha_W = \alpha/\sin^2\theta_W\) は追加パラメータなしで決定される。
	
	\subsubsection{強い結合}
	量子色力学は 8 個のグルーオン生成子を持つため、プランクスケールでの逆結合は
	\begin{equation}
		\alpha_s^{-1}(M_{\mathrm{Pl}}) = \left(\frac{3}{2}\right)^{\!8 + \beta\gamma\,S_{\mathrm{even}}/S_{\mathrm{odd}}}
		\approx 25.5 \qquad (\alpha_s \approx 0.039) .
	\end{equation}
	エネルギースケールの低下は YUCT では追加の調整殻の逐次活性化として記述され、各殻は逆結合に離散的な増分 \(\Delta L\) を寄与する。連続極限ではこれは QCD の対数的走りを再現するが、\(S_{\mathrm{odd}}\) と \(S_{\mathrm{even}}\) の整数倍の深度での自然な離散化を持つ。
	
	\subsubsection{湯川結合}
	荷電フェルミオン \(f\) の湯川結合は
	\begin{equation}
		y_f = \frac{\sqrt{2}\,m_f}{v},
		\qquad
		m_f = M_{\mathrm{Pl}}\left(\frac{2}{3}\right)^{\!96+k_f}\xi_f,
		\qquad
		v   = M_{\mathrm{Pl}}\left(\frac{2}{3}\right)^{\!96}\xi_W .
	\end{equation}
	したがって
	\begin{equation}
		y_f = \sqrt{2}\,\frac{\xi_f}{\xi_W}\,\left(\frac{2}{3}\right)^{\!k_f} .
	\end{equation}
	共鳴因子 \(\xi_f\) は、「YUCT 調整系統学」で導入された適合関数 \(C_{fH}(L_f,L_H)\) を通じて表現され、普遍的幅 \(\sigma \approx 0.20\) を持つ。したがってすべての湯川結合は、フェルミオンとヒッグス粒子の既知の調整深度によって固定される。
	
	\subsubsection{有効モードの離散梯子}
	図~\ref{fig:ladder} は、有効ゲージモード数 \(N_{\mathrm{eff}}\) が深度 \(L\)（エネルギー低下）とともにどのように変化するかを示している。各ステップは質量閾値に対応し、その閾値で粒子（または粒子族）が真空偏極への寄与を停止する。滑らかな包絡線は、観測される結合定数を決定する低エネルギー極限である。
	
	\begin{figure}[htbp]
		\centering
		\begin{tikzpicture}[scale=0.9]
			% 軸
			\draw[->] (0,0) -- (16,0) node[right] {\(L\)（深度）};
			\draw[->] (0,0) -- (0,8) node[above] {\(N_{\mathrm{eff}}\)};
			% 階段
			\draw[thick,blue] (0, 5.0) -- (1.5, 5.0);
			\draw[thick,blue] (1.5, 5.5) -- (3.2, 5.5);
			\draw[thick,blue] (3.2, 6.0) -- (5.0, 6.0);
			\draw[thick,blue] (5.0, 6.5) -- (7.0, 6.5);
			\draw[thick,blue] (7.0, 7.0) -- (9.0, 7.0);
			\draw[thick,blue] (9.0, 7.5) -- (11.0, 7.5);
			\draw[thick,blue] (11.0, 7.8) -- (13.0, 7.8);
			\draw[thick,blue] (13.0, 7.9) -- (15.0, 7.9) node[right,blue] {12.1333};
			% ラベル
			\node[below] at (1.5,0) {$L_W$ (96)};
			\node[below] at (3.2,0) {$L_Z$};
			\node[below] at (5.0,0) {$L_{\mathrm{top}}$};
			\node[below] at (7.0,0) {$L_{\mathrm{Higgs}}$};
			\node[below] at (9.0,0) {$L_{\tau}$};
			\node[below] at (11.0,0) {$L_{\mu}$};
			\node[below] at (13.0,0) {$L_e$ (107)};
			% 滑らかな包絡線
			\draw[red, dashed, thick] plot[smooth] coordinates {(0,5.0) (1.5,5.5) (3.2,6.0) (5.0,6.5) (7.0,7.0) (9.0,7.5) (11.,7.8) (13.,7.9)};
		\end{tikzpicture}
		\caption{有効ゲージモード数 \(N_{\mathrm{eff}}\) の調整深度 \(L\) に対する依存性。ステップは質量閾値に対応し、破線は微細構造定数を決定する低エネルギー極限である。}
		\label{fig:ladder}
	\end{figure}
	
	これらの結果により、標準模型のすべての無次元パラメータは、トポロジカル不変量 12、基本比 \(S_{\mathrm{odd}}/S_{\mathrm{even}}\)、普遍的修正 \(\beta\gamma\)、そして粒子の調整深度にまで遡ることができる。自由連続パラメータはもはや残らない。
	
	\subsubsection{標準模型における自由連続パラメータの完全な除去}
	
	上述の結果は、いったん YUCT フレームワークに埋め込まれると、標準量子場理論には \emph{自由連続パラメータがまったく残らない} ことを示している。表~\ref{tab:SMparams} は、標準模型の各独立な定数が、調整理論の基本数と関連粒子の調整深度によってどのように固定されるかをまとめている。
	
	\begin{table}[htbp]
		\centering
		\caption{YUCT における標準模型無次元パラメータの決定}
		\label{tab:SMparams}
		\small
		\begin{tabular}{@{}p{3.0cm} p{5.2cm} p{5.2cm}@{}}
			\toprule
			\textbf{パラメータ} & \textbf{YUCT における表現} & \textbf{主な参考文献} \\
			\midrule
			微細構造定数 \(\alpha\) &
			\(\alpha^{-1} = (S_{\mathrm{odd}}/S_{\mathrm{even}})^{12+\beta\gamma\,S_{\mathrm{even}}/S_{\mathrm{odd}}}\) &
			本文，付録 P，Ferra1.2 \\
			ワインバーグ角 \(\theta_W\) &
			\(m_W/m_Z\) から導出され，両方の質量は深度 \(L_W=96\) に従う &
			付録 \(\Lambda\) \\
			強い結合 \(\alpha_s\) &
			\(\alpha_s^{-1} = (S_{\mathrm{odd}}/S_{\mathrm{even}})^{8+\beta\gamma\,S_{\mathrm{even}}/S_{\mathrm{odd}}}\)（プランクスケール）；走りは離散深度殻から得られる &
			本文，付録 P \\
			湯川結合 \(y_f\) &
			\(y_f = \sqrt{2}\,(\xi_f/\xi_W)\,(S_{\mathrm{even}}/S_{\mathrm{odd}})^{k_f}\)，共鳴因子 \(\xi_f,\xi_W\) は適合行列から &
			付録 J，調整系統学 \\
			フェルミオン質量 & 半整数量子化 \(m_f = m_e\,q^{N_f}\)，\(q=(S_{\mathrm{odd}}/S_{\mathrm{even}})^{1/3}\) &
			付録 J，調整系統学 \\
			CP 破れ \(\theta\) パラメータ &
			\(\theta \sim (S_{\mathrm{even}}/S_{\mathrm{odd}})^{L_W/2 + S_{\mathrm{odd}}/S_{\mathrm{even}}}\) &
			「強い CP 問題の解決」 \\
			宇宙定数 \(\Omega_\Lambda\) &
			\(\Omega_\Lambda = 12/18 = 2/3\)（トポロジカル不変量） &
			本文，付録 A \\
			\bottomrule
		\end{tabular}
	\end{table}
	
	表のすべての項目は、調整ネットワークを特徴づける三つの整数（基準深度 \(L_W=96\)、ゲージモード数 \(N_{\mathrm{gauge}}=12\)、グルーオンモード数 \(N_{\mathrm{gluon}}=8\)）とともに、普遍比 \(S_{\mathrm{odd}}/S_{\mathrm{even}}\)、積 \(\beta\gamma\)、そしてフェルミオンとヒッグス粒子の離散調整深度 \(L_f\) のみに依存している。共鳴因子 \(\xi_f/\xi_W\) は調整可能なパラメータではなく、適合行列 \(C_{fH}\) から計算され、その幅 \(\sigma \approx 0.20\) は経験的に安定している（YUCT 調整系統学参照）。CP 奇 \(\theta\) パラメータも同様に、電弱スケールの深度と基本的な秩序–カオス偏移によって固定される。したがって、標準模型に現れる無次元定数の全体集合は、任意の連続入力なしに決定される。
	
	この構造的還元は、量子場理論を 19 個の自由パラメータを持つ理論から、その数値内容が YUCT の調整アーキテクチャによって完全に決定される有効記述へと変換するプログラムを完成させる。
	
	\section{統一調整場と二つのプロトコル}
	\label{sec:unified-protocols}
	
	上述の幾何学的張力フレームワークは、完全な 19 次元 YUCT 幾何学に埋め込むことでより正確に定式化できる。前に示したように、宇宙定数 \(\Omega_\Lambda=2/3\) は緊密ファイバー \(F_{18}\) のトポロジーから生じる。より深い分析は、\(\Theta_{\mu\nu}\) を生成する同じ調整場 \(\Psi_{MN}\) が二つの基本調整プロトコル（集中型 YPSDC と分散型 d‑YPSDC）も与えることを明らかにする（完全な導出は第一原理論文を参照）。統一作用は
	\begin{equation}
		S_{19} = \frac{1}{2\kappa_{19}} \int d^{19}X \sqrt{-G}\, R(G)
		+ \int d^{19}X \sqrt{-G}\,
		\Bigl[ -\frac14 \Tr(\Psi_{MN}\Psi^{MN})
		+ \frac12 G^{MN}\partial_M\phi \partial_N\phi
		- V(\phi) \Bigr],
		\label{eq:unified-gravity}
	\end{equation}
	ここで \(\phi = \ln(K_{\mathrm{eff}}/K_0)\) は本論文で使用される有効 4 次元スカラー場である。\(F_{18}\) でコンパクト化した後、\(\Psi_{MN}\) のゼロモードはちょうどスカラー \(\phi\) と式 (\ref{eq:Theta}) の幾何学的張力テンソル \(\Theta_{\mu\nu}\) を与える。
	
	二つのプロトコルは同じ場の異なる力学状態に対応し、索引 \(\kappa\) の起源によって区別される：
	\begin{enumerate}[leftmargin=*]
		\item \textbf{集中型 YPSDC。} 一つのエージェントが能動的に \(\Psi_{MN}\) を励起し、局所的な摂動を生成して他のエージェントに伝播させる。送信者は点状の源 \(D_M\Psi^{MN} = J^N_{\rm sender}\) として振る舞う。これはおなじみの通信の図であり、重力の文脈ではニュートンポテンシャルを生成する局所的な過密に対応する。
		\item \textbf{分散型 d‑YPSDC。} すべてのエージェントは \(\Psi_{MN}\) の長波長背景モードから同じ環境索引を受動的に読み取る。送信者は不要であり、相関は \(F_{18}\) の大域的トポロジーと共有辞書によって強制される。重力では、この状態は宇宙定数、暗黒エネルギー、および量子エンタングルメントで観測される非局所相関を生じさせる。
	\end{enumerate}
	
	状態間の遷移は無次元パラメータ \(\Theta = |J_{\rm agent}| / |J_{\rm env}|\) によって制御される。ここで \(J_{\rm agent}\) は能動的送信者の流れ、\(J_{\rm env}\) は宇宙論的背景の有効流である。\(\Theta\ll1\)（分散型状態）では場 \(\phi\) はほぼ一様で \(\partial_M\phi\approx0\) であり、幾何学的張力テンソルは有効宇宙定数と小さな暗黒物質成分に還元される。\(\Theta\gg1\)（集中型状態）では \(\phi\) の局所勾配が支配的となり、ニュートンポテンシャルと平坦な銀河回転曲線を説明する \(1/r\) 補正を生み出す。
	
	この統一は量子物理学に深遠な帰結をもたらす。量子エンタングルメントはもはや「不気味な遠隔作用」ではなく、単に分散型 d‑YPSDC の極限ケースである：エンタングルしたペアは共通の辞書（波動関数）を共有し、一方の粒子の測定は索引を提供し、場 \(\Psi_{MN}\) はその索引を同時に他方の粒子に届ける。ベル不等式の違反度は調整効率に関係し：
	\begin{equation}
		S = 2\sqrt{2}\,\frac{K_{\mathrm{eff}}}{K_{\mathrm{eff}}+1},
		\label{eq:bell-eff}
	\end{equation}
	したがって標準的な量子最大値 \(2\sqrt{2}\) は \(K_{\mathrm{eff}}\to\infty\)（完全な辞書）でのみ達成され、有限の効率では相関は弱くなる。
	
	このように、重力の幾何学的張力モデル、宇宙定数、フェルミオン質量階層性、そして量子非局所性はすべて、\(\mathcal{M}_{19}=M_4\times F_{18}\) 上の同じ実体 — 調整場 \(\Psi_{MN}\) — から導かれる。二つのプロトコルは別個の公理ではなく、同じ基礎物理学の導出された動力学的相である。二つの状態間の相転移とその観測的特徴の定量的分析は別の場所で提示される。
	
	\section{循環宇宙論}
	
	ポテンシャル \(V(\phi)\) は \(\phi=0\) に唯一の大域的極小を持つ。重力崩壊の間に、スカラー場は大きな正の値に駆動される。\(\phi\) が臨界閾値 \(\phi_{\mathrm{crit}}\) を超えると、有効圧力は十分に負になり、タキオン不安定性を引き起こす。条件は
	\[
	V''(\phi) > \frac{9}{4} H^2.
	\]
	Friedmann–Yakushev 方程式
	\[
	H^2 = \frac{8\pi G_0}{3} \left[ \rho_b + \frac{\kappa}{2}\dot\phi^2 + V(\phi) \right],
	\]
	を用い、バウンス近傍では \(\dot\phi^2\) が大きいことに注意すると、不安定性は次のときに起こることが分かる：
	\[
	e^{\lambda\phi_{\mathrm{crit}}} \approx \frac{9}{4} \frac{\kappa\dot\phi^2}{\lambda^2 V_0} + \cdots
	\]
	詳細な数値積分は \(\phi_{\mathrm{crit}} \approx \ln\Phi\) を示す。ここで \(\Phi \approx 1.618\) は黄金比である。その後、場は急速に \(\phi=0\) へと転がり落ち、蓄えられたエントロピーを放出し、短い期間の指数関数的膨張（インフレーション）を駆動する。このサイクルは無限に繰り返され、初期特異点を回避する。
	
	\section{宇宙の大規模構造と基本スケールの本質}
	\label{sec:sponge}
	
	調整場 \(\phi\) の方程式は局所重力を決定するだけでなく、宇宙の全体的なアーキテクチャも形作る。このセクションでは、「スポンジ状」の宇宙網、因果的に切り離された子宇宙の可能性、そしてプランク長の非普遍性が、どのように幾何学的張力フレームワークの自然な結果として現れるかを示す。
	
	\subsection{宇宙網のスポンジ状形態}
	物質優勢時代、場方程式~\eqref{eq:phieom} は異なる調整深度 \(L\) の領域を分離する定常的な非線形解を許容する。二つの領域 \(\phi_1\) と \(\phi_2\) の間の 1 次元平面壁のプロファイルは
	\begin{equation}
		\phi(x) = \frac{2}{\lambda}\,\ln\!\left[
		\frac{1 + \tanh(\lambda\sqrt{V_0/\kappa}\, x/2)}
		{1 - \tanh(\lambda\sqrt{V_0/\kappa}\, x/2)}
		\right],
		\label{eq:domain-wall}
	\end{equation}
	であり、二つの真空の間を滑らかに補間する。そのような壁の表面張力は \(\sigma \simeq \sqrt{\kappa V_0}/\lambda^2\) である。
	
	三次元では、これらの壁は連結した細胞ネットワークを形成する：節点は深いクラスタ（\(\phi\) 大）、フィラメントは壁の交差、ボイドは \(\phi \approx 0\) の領域である。これはまさに銀河サーベイで観測される「宇宙スポンジ」である。壁の重力レンズ信号は \(\phi\) の勾配によって増強され、有効暗黒物質面密度を与える：
	\begin{equation}
		\Sigma_{\rm DM}^{\rm eff} = \frac{\kappa}{4\pi G_0} \int_{-\infty}^{+\infty} \!dz\; \nabla^2\delta\phi(z)
		\simeq \frac{\sqrt{\kappa V_0}}{2\pi G_0\,\lambda},
		\label{eq:wall-lensing}
	\end{equation}
	これはスタックされた弱レンズデータで検証可能である。
	
	\subsection{切り離された子宇宙と「不定距離」}
	\(\phi\) の局所的な揺らぎが臨界値 \(\phi_{\rm crit}=\ln\Phi\) を超えると、\textbf{局所的な} 相転移が起こる：新しい真空の小さな泡が核形成し、膨張し、因果的に切り離された領域 — 子宇宙 — を生成する。二つの宇宙は非常に急峻な \(\nabla\phi\) の表層を持つドメイン壁によって隔てられている。壁を横切る曲線のアフィンパラメータは発散するため、古典的な測地線によって二つの間の距離を測定することはできない。唯一意味のある「距離」は \textbf{調整深度差}
	\begin{equation}
		\Delta L = L_{\rm parent} - L_{\rm child},
		\label{eq:deltaL}
	\end{equation}
	であり、これはトポロジカル量子数である。したがって、宇宙間の「不定距離」は調整層の離散構造の直接的な結果である。
	
	\subsection{局所的で環境依存的なスケールとしてのプランク長}
	YUCT では、基本長は深度 \(L\) を持つ最も軽い調整クラスタのコンプトン波長である。質量公式 \(m \simeq M_{\rm Pl}\,(2/3)^L\) を用いると、この長さは
	\begin{equation}
		\ell_{\rm coord}(L) = \frac{\hbar}{m c}
		= \frac{\hbar}{M_{\rm Pl}c} \left(\frac{3}{2}\right)^{\!L}
		= \ell_{\rm Pl} \left(\frac{3}{2}\right)^{\!L}.
		\label{eq:length-from-L}
	\end{equation}
	現在の時代の真空に対しては \(L_{\rm vacuum}\to\infty\) であるが、有限のクラスタはスケール \(\ell_{\rm coord}(L) \gg \ell_{\rm Pl}\) を探る。通常のプランク長は、この一般公式の \(L=1\) という特別な場合、すなわち今日の低エネルギー現象を支配する調整層を特徴づけるものに過ぎない。
	
	したがって、重力定数も層に依存する：
	\begin{equation}
		G_{\rm eff}(\phi) = \frac{G_0}{1-4\pi G_0\alpha_2\phi^2},
		\label{eq:Geff-layer}
	\end{equation}
	そのため \(\phi\) が大きい領域では重力が弱くなる。これはハッブル・テンションの自然な解決を提供する：膨張率の局所測定は、CMB とはわずかに異なる \(\phi\) 環境を探り、その結果として異なる \(H_0\) の推定値をもたらす。この効果の定量的研究は別の場所で提示される。
	
	\section{検証可能な予測}
	
	\subsection{銀河回転曲線}
	修正重力項は、宇宙暗黒物質分率によって規格化された平坦な回転曲線を予測する。個々の銀河への詳細なフィッティングは別の研究で示される。
	
	\subsection{宇宙構造の成長}
	調整場は線形成長因子を修正する。標準的な摂動論を用いると、赤方偏移 \(z=0.5\) における \(\Lambda\)CDM からの分数偏差は約 \(2\)–\(3\%\) であり、Euclid や DESI によって測定可能である。
	
	\subsection{宇宙論パラメータの制約}
	このモデルは、\(\Omega_\Lambda = 2/3\) および \(\Omega_{\mathrm{DM}}^{\mathrm{eff}} = 1 - 0.05 - 2/3 \approx 0.2833\) を正確に予測する。これらの値は Planck、DES、LSST データと対照できる。
	
	\section{結論}
	我々は、YUCT の情報論的公理から重力を導出した。幾何学的張力テンソル \(\Theta_{\mu\nu}\) は暗黒物質と暗黒エネルギーを統一し、宇宙定数はトポロジーによって固定され、循環宇宙史が自然に現れる。すべてのパラメータはプランクスケールと辞書定数で表現される。この理論は現在の宇宙論サーベイで反証可能である。
	
	\begin{thebibliography}{99}
		\bibitem{Y-appA} A. V. Yakushev, YUCT Lagrangian V35.0 (付録 A).
		\bibitem{Y-appL} A. V. Yakushev, Fractal Coordination Error Scaling (付録 L).
		\bibitem{Y-appM} A. V. Yakushev, Cosmology -- Inflation as a Coordination Phase Transition (付録 M).
		\bibitem{Y-appG} A. V. Yakushev, Resolution of the Quantum Gravity Problem (付録 G).
	\end{thebibliography}
	
\end{document}